\paragraph{Betriebsverstärkung}
Die experimentelle Messung liefert deutlich niedrigere Fehler als die theoretische Berechnung. Dies liegt zum einen am großen Widerstandsfehler von \(\qty{10}{\percent}\), zum anderen daran, dass ein linearer Fit nur Fehler im \verb|y-data| Set berücksichtigt, Fehler in der Eingangsspannung werden also nicht miteinbezogen (ist aber wsl. vernachlässigbar im Gegnesatz zum hohen relativen Fehler der Widerstände). Durch diese hohen theoretischen Fehler liegt die Sigma-Abweichung der zwei Methoden ziemlich niedrig, siehe \cref{table:Vergleich_Vexp}. Betrachtet man als Ersatz-Parameter die absolute Abweichung, so steigt diese mit zunehmenden Rückkopplungswiderstand. Bei Wechselspannungen liegt dies vielleicht an dem beim Frequenzganganalyse (\cref{fig:Nr.Frequenzgang.png}) beobachteten Effekt, dass die Betriebsverstärkung ihre Konstanz für höhere \(R_g\) früher verliert und dadurch die Abschätzungen in \cref{eq:Abschätzungen} nicht mehr gelten, sodass sich die Diskrepanz zwischen vereinfachter Formel und Realität in den abweichenden Ergebnissen niederschlägt.\\
Bei Gleichspannungsmessungen weiß ich nicht, welche Gründe die steigende absolute Abweichung bei steigendem \(R_g\) hat.

\paragraph{Frequenzgang}
Warum die niedrigere Abfallsgrenzfrequenz (dass für höhere \(R_g\) die Betriebsverstärkung früher in die Leerlaufverstärkung übergeht) liegt meiner Meinung nach daran, dass gemäß \cref{eq:Abschätzungen} für höhere \(R_g\) \(\frac{R_e}{R_g}\) nicht mehr groß gegen \(\frac{1}{V_0}\) ist, da \(V_0\) gemäß Skript\autocite{Skript_2.2} für hohe Frequenzen sinkt.\\
Ansonsten wurden alle weiteren Dinge schon in der Auswertung hervorgehoben. 

\paragraph{Recheckpulsform}
Die Frage, warum die unscharfen Kanten nur bei der linken Flanke auftreten, ist weiterhin unklar (hast du ne Idee?). Ansonsten sieht man sehr schön die Eigenschaften der verschiedenen Schaltungen. Zusammenfassend, wie in der Auswertung schon beschrieben, ist die Schaltung mit der kleinsten Veränderung der Rechteckform diejenige, die weder tiefe Frequenzen, noch hohe Frequenzen filtert, also weder einen Eingangshochpass (Schalter S1 auf Stellung 2) noch Ausgangshochpass, sondern geringsten Rückkopplungswiderstand (Schalter S2 auf Stellung 3) besitzt. Wichtig zu bemerken: Der Kompromiss ist jedoch eine niedrigere Pulshöhe (Maximalspannung), da die Betriebsverstärkung durch geringen Rückkopplungswiderstand kleiner ist. Das ist in den Oszilloskopbildern grafisch nicht zu sehen, da diese alle auf dieselbe vertikale Höhe skaliert wurden. Glücklicherweise berechnet das Oszi jedoch automatisch den Spitze-Spitze-Spannungswert \(V_{SS}\), dieser beträgt bei \cref{fig:TEK0002.JPG} (beste Erhaltung der Rechteckform) \(V_{SS}=\qty{1.54}{\volt}\), bei \cref{fig:TEK0004.JPG} (\(R_g=\qty{680}{\kilo\ohm}\) maximal, Schalter S2 auf Stellung 1, S1 weiterhin auf 2) ist \(V_{SS}=\qty{20.8}{\volt}\). 

\subsection{Fazit}  
Netter Versuch, die Ergebnisse sind ziemlich verständlich, die Auswertung erzeugt nachvollziehbare Beobachtungen. Aber irgendwie wäre eine bessere Einführung in die Funktionsweise eines invertierenden und nicht invertierenden Operationsverstärker ganz hilfreich und interessant.
